\begin{theorem}[A uniform boundary-column action passes to the actual cut]%
    \label{thm:peps_physical_cut_column_action}
    \lean{TNLean.PEPS.mul_regularPhysicalCutMatrix_eq_smul_of_openRegion_eigen,
        TNLean.PEPS.mul_regularPhysicalCutMatrix_eq_ite_of_openRegion_eigen}
    \leanok
    \uses{thm:peps_regular_open_region_projector}
    Let $B_R$ be the actual open-region coefficient matrix of a finite regular
    tensor network, and let $K_R=B_R B_{R^c}^{\mathsf T}$ be its physical cut
    matrix. If a physical matrix $Q$ satisfies
    $Q\Psi_R(\theta)=z\Psi_R(\theta)$ for every crossing configuration $\theta$,
    then $QK_R=zK_R$. In particular, if one condition $P$ selects every
    actual group-labelled column or annihilates every such column, the same
    condition selects $K_R$ or annihilates it.
    This is the contraction step used in the local measurements of
    \cite[Theorem~6.15]{Schuch2010PEPS}.
\end{theorem}

\begin{proof}\leanok
    The column identities give $QB_R=zB_R$. Therefore
    $QK_R=(QB_R)B_{R^c}^{\mathsf T}=zK_R$.
    For the conditional statement, reindex the finite group labels and take
    $z=\mathbf1_P$.
\end{proof}

\begin{definition}[The outer walk and routed joint-flux assignment]%
    \label{def:peps_torus_joint_flux_geometry}
    \lean{TNLean.PEPS.torusTwoPlaquetteOuterWalk,
        TNLean.PEPS.torusTwoPlaquetteOuterWalk_support,
        TNLean.PEPS.torusJointFluxAssignment,
        TNLean.PEPS.regularDirectedTransport_torusJointFluxAssignment_up}
    \leanok
    \uses{thm:peps_translated_two_plaquette_geometry}
    On a torus of horizontal period at least four and vertical period at least
    three, let $R_v$ be the six-site block of two adjacent plaquettes based at
    $v=(x,y)$. Its counterclockwise outer walk makes two rightward steps, one
    upward step, two leftward steps, and one downward step. Its visited sites
    are exactly $R_v$.

    For $a,b\in G$, define the routed assignment $u_{a,b}$ by giving upward
    transport $(ab)^{-1}$ to the left vertical bond and $b^{-1}$ to the middle
    vertical bond. Every other transport is the identity. Each ordered bond
    receives the specified element when its ordering agrees with the upward
    direction, and its inverse otherwise. Thus this definition also applies
    at the vertical seam. This is the routed geometry of the joint flux
    measurement in \cite[braiding passage]{Schuch2010PEPS}.
\end{definition}

\begin{theorem}[The actual individual and joint holonomies]%
    \label{thm:peps_torus_joint_flux_holonomy}
    \lean{TNLean.PEPS.regularWalkHolonomy_torusJointFluxAssignment}
    \leanok
    \uses{def:peps_torus_joint_flux_geometry,
        def:peps_torus_plaquette_flux_geometry}
    The counterclockwise left, right, and outer walks have respective
    holonomies $a$, $b$, and $ab$. The formulas hold at every translated
    position, including both periodic seams.
\end{theorem}

\begin{proof}\leanok
    In the reverse-ordered holonomy convention the left loop gives
    $(ab)b^{-1}=a$, the right loop gives $b$, and the outer loop gives $ab$.
    All horizontal transports and the rightmost vertical transport are the
    identity. The ordered-edge inversion preserves these directed values
    at the seams.
\end{proof}

\begin{theorem}[One physical measurement detects the joint flux]%
    \label{thm:peps_torus_joint_flux_physical_measurement}
    \lean{TNLean.PEPS.IsGIsometric.exists_torusJointFlux_physicalMeasurement}
    \leanok
    \uses{thm:peps_translated_two_plaquette_geometry,
        thm:peps_torus_joint_flux_holonomy,
        thm:peps_closed_walk_class_physical_measurement,
        thm:peps_physical_cut_column_action}
    Let the original four-leg tensor be $G$-isometric for the regular action.
    Under the stated period bounds, for every position $v$ there is a complete
    projective measurement $(Q_C)_C$ on the six original spins of $R_v$.
    Its outcomes are the conjugacy classes of $G$ and one additional outcome
    for the unused physical subspace. Each $Q_C$ is positive and Hermitian,
    $\sum_C Q_C=I$, and $Q_CQ_D=\delta_{C,D}Q_C$.

    This one measurement is chosen before $a,b$, every crossing configuration
    $\theta$, and every common exterior and crossing assignment $u$. For the
    actual routed open column and actual global cut matrix, it satisfies
    \begin{align}
        Q_C\Psi_{R_v}(u_{a,b};u,\theta)
         & =\mathbf 1_{C=C[ab]}\Psi_{R_v}(u_{a,b};u,\theta), \\
        Q_CK_{R_v}(u_{a,b};u)
         & =\mathbf 1_{C=C[ab]}K_{R_v}(u_{a,b};u).
    \end{align}
    Here the internal bonds carry $u_{a,b}$, whereas all exterior and crossing
    bonds retain $u$. This proves the regular finite-torus joint measurement
    step described in \cite[braiding passage]{Schuch2010PEPS}.
    No braiding transformation, energy, or parent-Hamiltonian membership is
    asserted; the broader source scope remains separate as recorded in
    \cite{gap:rmp_peps_quantum_double_g_isometry}.
\end{theorem}

\begin{proof}\leanok
    The outer walk induces a closed walk in the actual connected six-site
    block. Native four-leg $G$-isometry gives incident-edge $G$-isometry at each
    of its sites. Apply the closed-walk class measurement to this walk.
    The proved outer holonomy is $ab$, independently of all exterior and
    crossing assignments. The column identity then passes to the actual
    cut matrix by the uniform boundary-column theorem.
\end{proof}
